\begingroup
\renewcommand{\arraystretch}{1.45}
\fancyhead[L]{\footnotesize CONCERNING THE THEORY OF THE NEW TRANSCENDENTAL FUNCTIONS. V.}

\Needspace{22\baselineskip}
\par\medskip\Needspace{5\baselineskip}\pdfbookmark[1]{Art. 16}{transscendentes-v-art-16}\centerline{[16.]}\nobreak\smallskip
If throughout the interior of a bounded figure
\[
\frac{\mathrm{dd}V}{\mathrm{d}x^2}
+\frac{\mathrm{dd}V}{\mathrm{d}y^2}=0
\]
while on the boundary $V$ is constant $=A$, then necessarily $V=A$ throughout the region as well


\Needspace{20\baselineskip}
Proof. Let $\mathrm{d}M$ be an element of the area and
\[
\begin{gathered}\int\Bigl\{
\Bigl(\frac{\mathrm{d}V}{\mathrm{d}x}\Bigr)^2
+\Bigl(\frac{\mathrm{d}V}{\mathrm{d}y}\Bigr)^2
\\[8pt]\quad{}+\Bigl(\frac{\mathrm{dd}V}{\mathrm{d}x^2}
+\frac{\mathrm{dd}V}{\mathrm{d}y^2}\Bigr)(V-A)
\Bigr\}\mathrm{d}M=\Omega\end{gathered}
\]


\Needspace{25\baselineskip}
\SourceLeaf{489}

If $V$ were not constant, the integral would obviously be \emph{positive}. But the integral is also
\[
\int(V-A)\left(\frac{\mathrm{d}V}{\mathrm{d}x}\mathrm{d}y
-\frac{\mathrm{d}V}{\mathrm{d}y}\mathrm{d}x\right)
\]
taken around the perimeter of the figure, hence $=0$. Since this contradicts the preceding result,
the supposition is inadmissible.

The same applies with three indeterminates.

Various beautiful consequences can be drawn from this.

A point cannot be in stable equilibrium inside a hollow space unless no action occurs anywhere within the entire space,
either in the case of repulsion or in that of attraction.


\Needspace{20\baselineskip}
\par\medskip\Needspace{5\baselineskip}\pdfbookmark[1]{Art. 17}{transscendentes-v-art-17}\centerline{[17.]}\nobreak\smallskip
That there can be no \emph{negative} parts on the periphery of an equilibrium figure
is proved as follows. Suppose $AB$ were negative; then $V-A$ ($A$ being the value
of $V$ on the periphery) would be positive just outside $AB$, and negative beside the other parts. There will therefore be a finite region in which this positive value holds;
let it lie within $ABCDE$; but then, since it is $0$ on the boundary, it would be constant
throughout this region, which is a contradiction.

The same reasoning can be applied in three dimensions.


\Needspace{18\baselineskip}
In three dimensions, for the surface $s$ on which $V=\mathrm{Const.}$, there is also the
beautiful theorem that $\int\frac{\mathrm{d}V}{\mathrm{d}\rho}\mathrm{d}s$,
when $\rho$ is normal to $s$, $=4\pi M$, where $M$ is
the entire mass (presumably always the mass enclosed by such a surface,
thus excluding the mass outside it).

Presumably the theorem holds for \emph{every} enclosing surface, without $V$
needing to be constant.


\Needspace{20\baselineskip}
Let there be two surfaces, one enclosing the other, on which $V$ is constant,
namely $V=A$, $V=B$, respectively; let $\mathrm{d}W$ be an element of the space between them,
and $p$ the force of attraction at each point. Then
\[
\int pp\,\mathrm{d}W=(B-A)\, .\,4\pi M
\]
where $M$ is the mass in the inner region, it being assumed that there is no attracting mass
in the region $W$.

\par\nopagebreak\begin{center}\rule{93.96bp}{0.35pt}\end{center}

\clearpage
\endgroup
